\section{The Objective}
\label{sec:v8-objective}
\label{sec:v8-theorem-ladder}

The C-TreePO objective is two-channel: one channel fits the
document-level expert target, and one channel pushes \(g\) toward
the local laws. We write the objective, then climb the theorem
ladder (preservation, optimization, certification) to read off what
the laws buy.

\subsection{The C-TreePO Objective}

Section~\ref{sec:v8-ctree-math} gives the structural goal: learn
one state map \(g\) whose node states can replace the document
spans they represent for oracle-compatible readouts. We now write
the loss that pushes \(g\) toward that goal. The shape is
two-channel. The root channel asks the predicted document score to
match the expert target. The local channel asks every realized state
to behave like a valid state in its own right.

Let \(L_{\mathrm{root}}(g)\) be the root-level loss, measured
against the available expert target. In the Manifesto LLM
replication, \(L_{\mathrm{root}}\) measures the discrepancy between
the predicted root score and the expert-survey mean; the precise
per-node agreement metric used by the prompt-and-demonstration
optimizer is given in Section~\ref{sec:v8-benoit-tree-setup}.

For the local channel, attach one node loss to each \(v\in V(T)\).
The form of the loss depends on which law applies at \(v\). At a
leaf, the loss measures C1: does the leaf summary preserve the
evidence in the raw text? At a state passed back through \(g\), the
loss measures C2: does summarizing again give the same oracle-facing
answer? At an internal merge, the loss measures C3: does merging
the two children preserve the joint evidence of the parent span?
If every check could be measured by the target oracle, the local-law
loss at node \(v\) would be
\[
  \ell_v^\ast
  =
  \ell_{\mathrm{law}}(v;f^\ast,g).
\]
The star marks the target quantity: the local-law loss as it would be
measured by the expert oracle, not by the cheaper proxy used during
dense training.

Let \(d(v)\) be the depth of node \(v\), with the root at depth \(0\),
and let \(\gamma\in[0,1]\) be a depth-discount weight. Each node enters
the local-law sum with weight \(\gamma^{d(v)}\), so the root counts
fully and deeper nodes count progressively less. \(\gamma=1\) treats
every node equally; smaller \(\gamma\) emphasizes nodes near the root
and damps deeper nodes. This is the same discount pattern used in
discounted-reward objectives. We report \(\gamma\) as a declared knob
even when \(\gamma=1\) is the operational default, so deeper-node
weighting is an explicit design choice rather than a hidden default.
The oracle local-law channel is
\begin{equation}
\label{eq:v8-local-law-estimand}
  L_{\mathrm{law}}^\ast(g,T)
  =
  \sum_{v\in V(T)}\gamma^{d(v)}\ell_v^\ast .
\end{equation}
The theorem-facing objective is
\begin{equation}
\label{eq:v8-main-objective}
  J_\lambda^\ast(g)
  =
  (1-\lambda)L_{\mathrm{root}}(g)
  +
  \lambda L_{\mathrm{law}}^\ast(g,T),
  \qquad 0\le\lambda\le 1 .
\end{equation}
Here \(\lambda\) is the analyst-chosen local-law share. \(\lambda=0\)
trains on the document-level expert target alone, with no local-law
pressure on \(g\). \(\lambda=1\) trains on local laws alone and lets
the root drift. Intermediate values balance the two. The root term
names the document score we ultimately report; the local term names
the state-validity pressure that makes the theorem relevant to the
fitted tree.

\subsection{The Theorem Ladder}

The ladder has three levels. Preservation is the structural level:
exact C1/C2/C3 laws let node and root states replace the raw spans
they represent. Optimization adds a downstream compatibility
premise: the loss or preference objective must depend on the input
only through the preserved oracle interface. Certification relaxes
exact validity to an estimated distortion budget on the realized
tree.

\begin{table}[t]
\centering
\small
\begin{tabular}{@{}p{0.24\linewidth}p{0.33\linewidth}p{0.34\linewidth}@{}}
\toprule
Stack & Added premise & Application-facing conclusion \\
\midrule
Preservation & C1/C2/C3 hold on the realized C-Tree for \(f^\ast\) &
Node and root states can replace their raw document spans for
oracle-compatible readouts. \\
Optimization & Preferences or losses depend on the input only through the
preserved oracle interface &
Training on valid root summaries has the same population target as training
on full documents. \\
Certificate & Approximate local-law distortion is bounded or estimated from
sampled nodes with logged propensities &
Approximate preservation becomes a finite-sample bound on downstream
objective drift. \\
\bottomrule
\end{tabular}
\caption{\textbf{The theorem ladder.} Each tier keeps the task target
explicit: the oracle is the measurement the analyst wants to preserve.
Appendix~\ref{app:v8-theorem-certificate-details} records the
assumption-to-claim crosswalk for each result.}
\label{tab:v8-theorem-ladder}
\end{table}

The preservation stack is Theorem~\ref{thm:v8-root-preservation} and
Corollary~\ref{cor:v8-preferences-pass-to-root}. It is the
local-to-global claim: if every realized state remains in the same
\(f^\ast\)-fiber as the span it represents, then the root state
remains in the same fiber as the whole document.

The optimization stack adds an oracle-compatibility premise for the
downstream training object. The Manifesto LLM replication uses a
regression target, where this premise holds automatically:
\(L_{\mathrm{root}}\) sees the document only through its predicted
score. The premise becomes load-bearing when the framework is
applied to preference-learning or RL-style downstream tasks. For
pairwise preference learning, a comparison rule is compatible when
replacing a document by an equivalent root state does not change
the preferred answer. This covers DPO-style pairwise objectives
\citep{BradleyTerry1952,RafailovEtAl2023DPO} and GRPO-style group
objectives only when the prompt, response generator, policy,
ranker, or reward channel depends on the document through the
preserved oracle. A preference for formatting, rhetoric, or style
outside the chosen oracle interface needs a wider oracle.

\begin{proposition}[Population Preference Objective Equivalence]
\label{prop:v8-manifesto-preference-equivalence}
Fix a population of documents and a realized C-Tree construction for each
document. If C1/C2/C3 hold exactly for \(f^\ast\), and if the downstream
population integrand (including any pair generator, response generator,
policy, ranker, or reward channel used by the method) is oracle-compatible in
the input argument, then replacing each document by its C-Tree root state
leaves the population preference objective unchanged.
\end{proposition}

\begin{proof}[Proof sketch]
Assume C1, C2, and C3 hold exactly across the document population
and that the downstream population integrand is oracle-compatible
in its input argument. Theorem~\ref{thm:v8-root-preservation} gives
the per-tree root statement \(z_{\mathrm{root}}\sim_\ast x\) for
every document. Oracle-compatibility says that replacing an input
by another input in the same oracle fiber leaves the whole
integrand unchanged, including any generated comparisons,
responses, rewards, or ranks that the method averages over. Apply
that statement to each document's root state and raw text. The
integrand coincides at every realization, and expectation over the
document population preserves the identity. The full argument is in
Appendix~\ref{app:v8-full-proofs}.
\end{proof}

The proposition says that root summaries can serve as training
inputs for preference tasks whose input dependence is already
captured by the chosen oracle. Exact local laws handle the
compression step; the oracle-compatibility premise handles the
downstream objective.

\subsection{Approximate Preservation}

Exact local laws are the clean theorem endpoint. The deployed LLM
setting is approximate. The realized state at each node differs from
the raw span by some amount, and the realized score at the root
differs from the expert mean by some amount. Three quantities organize
that approximation.

\(G_{\mathrm{meth}}\) is the drift in whatever downstream quantity the
analyst is reporting (a population correlation, a population MSE, or a
preference-task population gap). This is the error the analyst cares
about.

\(\Delta_R\) is the average distortion between a produced state and
the raw span it represents, measured by the oracle. The subscript
\(R\) marks the represented-span comparison.

\(C_{\mathrm{meth}}\) is a method-specific transport constant. It
captures how much the downstream quantity can move when the
representation moves: a small \(C_{\mathrm{meth}}\) means a smooth
downstream method that tolerates small representation distortion well,
and a large \(C_{\mathrm{meth}}\) means a sharper method where small
representation errors translate into bigger reported drift.

\begin{proposition}[Population Gap Bound]\label{prop:v10-population-gap}
Fix a downstream method with transport constant \(C_{\mathrm{meth}}\)
and represented-span distortion \(\Delta_R\). The drift in the
reported quantity is bounded by
\begin{equation}
\label{eq:v8-population-gap}
  |G_{\mathrm{meth}}|
  \le
  C_{\mathrm{meth}}\,\Delta_R .
\end{equation}
\end{proposition}

\begin{proof}[Proof sketch]
The premise that defines \(C_{\mathrm{meth}}\) is the Lipschitz
inequality: when the representation moves by \(\Delta\) in the oracle
metric, the downstream quantity moves by at most
\(C_{\mathrm{meth}}\Delta\) at any single document. Apply this
pointwise inequality at every document in the population. The
triangle inequality gives
\(|\mathbb E[\mathrm{drift}]|\le\mathbb E[|\mathrm{drift}|]\), and
the Lipschitz premise bounds the last expectation by
\(C_{\mathrm{meth}}\mathbb E[d_Y(f^\ast(z),f^\ast(x))]\). The final
expectation is \(\Delta_R\).
The full argument, with the regression, DPO, and GRPO instances
stating the form of \(C_{\mathrm{meth}}\) explicitly, is in
Appendix~\ref{app:v8-full-proofs}.
\end{proof}

For a regression case (such as the Manifesto LLM replication),
\(C_{\mathrm{meth}}\) is the Lipschitz constant of the score
readout: how far the predicted score can move when the input state
moves a little. For DPO, \(C_{\mathrm{meth}}\) packages a
policy-Lipschitz condition on the log-probability ratios. For
group-ranking or RL-style objectives, \(C_{\mathrm{meth}}\) packages
the ranker, reward, clipping, and KL terms. The constant belongs to
the downstream training method; the C-Tree audit supplies
\(\Delta_R\).

What's left is to make \(\Delta_R\) operational: sample realized
tree nodes, log inclusion probabilities, and turn sampled distortion
into a finite-sample bound on
\(|G_{\mathrm{meth}}|\).\footnote{Appendix~\ref{app:v8-neural-operator-route}
gives the function-class route to \(\Delta_R\) through approximation
properties of \(g\).} The next section does that.
Appendix~\ref{app:v8-full-proofs} gives the proof details for the
preservation, preference-transport, and certificate statements used
in this ladder.
